\chapter{Збір даних.}

Першим етапом практично будь-якого проєкту з інтелектуального аналізу даних, машинного навчання 
або аналітики є збір даних (Data Collection). Якість отриманих результатів значною мірою залежить 
від якості та повноти вихідних даних. Навіть найскладніші алгоритми машинного навчання не здатні 
компенсувати відсутність важливої інформації або низьку якість даних, тому процес їх отримання є 
одним із найважливіших етапів підготовки аналітичного проєкту.

Сучасні дані можуть надходити з великої кількості різноманітних джерел. Це можуть бути 
реляційні бази даних, електронні таблиці, CSV-файли, журнали подій (logs), 
сенсори Інтернету речей (IoT), соціальні мережі, веб-сайти, а також численні 
програмні інтерфейси (API), що надають доступ до інформаційних сервісів. У 
деяких випадках дані вже доступні у вигляді готових файлів і можуть бути 
безпосередньо завантажені для подальшого аналізу. В інших ситуаціях 
необхідно виконувати запити до баз даних або зовнішніх інформаційних 
систем для отримання актуальної інформації. 

Автоматизований збір даних є значно швидшим, дозволяє обробляти великі обсяги інформації, 
зменшує кількість людських помилок і забезпечує регулярне оновлення даних без постійного 
втручання користувача. Основними труднощами під час збору даних є різні формати даних, 
неоднакова структура, дублікати, пропущені значення, обмеження доступу, необхідність 
автентифікації та різна якість отриманої інформації.

В цій темі ми поговоримо про збір даних з двох основних джерел:

\begin{itemize}
\item запит до API (більшість з яких сьогодні є веб-орієнтованими); та

\item вилучення даних з веб-сторінки за допомогою технологій веб-скрапінгу (Web Scraping).
\end{itemize}



\section{Збір даних з веб-джерел.}

Веб-джерело даних — це будь-який ресурс у мережі Інтернет, з якого можна отримати інформацію, 
наприклад веб-сторінка, REST API або онлайн-база даних. Для збору даних можуть використовуватися 
веб-сайти, REST API, XML- та JSON-сервіси, відкриті набори даних, хмарні сервіси та веб-додатки. 

Сьогодні переважна більшість автоматизованих систем обміну даними використовує протокол HTTP 
(HyperText Transfer Protocol). Хоча спочатку HTTP був створений для передачі веб-сторінок між 
браузерами та веб-серверами, з часом він став універсальним механізмом взаємодії між 
програмними системами. Саме на основі HTTP функціонує більшість сучасних веб-сервісів, 
REST API та хмарних платформ.

Під час збору даних клієнтська програма надсилає HTTP-запит до віддаленого сервера, а 
сервер повертає відповідь, яка може містити HTML-документ, JSON-структуру, XML-файл, 
зображення або інші типи даних. Розуміння принципів роботи HTTP є необхідною передумовою 
для ефективного використання веб-джерел даних, тому в наступних розділах буде розглянуто 
основні типи HTTP-запитів, структуру відповідей сервера та способи обробки отриманих даних у Python.

HTTP-відповідь — це повідомлення, яке сервер надсилає клієнту після обробки HTTP-запиту. 
Вона містить код стану, заголовки та тіло відповіді, у якому можуть передаватися HTML, JSON, 
XML або інші дані. 

Веб-скрапінг — це автоматизований процес отримання та вилучення інформації з HTML-сторінок 
веб-сайтів для її подальшого аналізу або обробки. Процес веб-скрапінгу включає отримання 
HTML-сторінки, аналіз її структури, пошук потрібних елементів, вилучення необхідних даних 
та їх подальше збереження або обробку.

Процес веб-скрапінгу складається із надсилання запиту, одержання відповіді, парсингу тексту, 
витягу очищення та збереження даних.

\begin{lstlisting}
import requests
response = requests.get("https://fmi.chnu.edu.ua/")

print("Status Code:", response.status_code)
print("Headers:", response.headers)

print(response.text)
\end{lstlisting}

Результат:

\begin{console}
Status Code: 200
Headers: {'Date': 'Fri, 12 Sep 2025 05:44:43 GMT', 'Content-Type': 'text/html; charset=utf-8' ...
\end{console}

Коди стану HTTP повідомляють клієнту про результат обробки запиту сервером. Код 200 (OK) 
означає, що запит успішно виконано. Код 404 (Not Found) вказує на те, що запитуваний ресурс 
не знайдено. Код 500 (Internal Server Error) свідчить про виникнення внутрішньої помилки на 
сервері під час обробки запиту.

HTTP GET є найпоширенішим методом, але існують також методи PUT, POST, DELETE, які 
змінюють деякі параметри на сервері.

Параметри URL можна вказати за допомогою бібліотеки запитів, як показано нижче:

\begin{lstlisting}
params = {"sa":"t", "rct":"j", "q":"", "esrc":"s",
"source":"web", "cd":"9", "cad":"rja", "uact":"8"}
response = requests.get("http://www.google.com/url", params=params)

print("Final URL:", response.url)
print(response.text)
\end{lstlisting}

Результат:

\begin{console}
Final URL: http://www.google.com/url?sa=t&rct=j&q=&esrc=s&source=web&cd=9&cad=rja&uact=8
<html lang="uk"><head><meta http-equiv="Content-Type" content="text/html; charset=ISO-8859-1"><title>&#1055 ...
\end{console}


\section{RESTful API}

Якщо ви перейдете від простого запиту веб-сторінок до веб-API, ви, швидше за все, зіткнетеся з REST 
API (Representational State Transfer). REST — це скоріше архітектура дизайну, але ось кілька 
ключових моментів:

\begin{itemize}
\item Використовує стандартний інтерфейс і методи HTTP (GET, PUT, POST, DELETE)
\item Stateless – сервер не запам'ятовує, що ви робили
\item Практичне правило: якщо ви надсилаєте ключ свого облікового запису разом із 
кожним викликом API, ви, ймовірно, використовуєте REST API
\end{itemize}

REST API забезпечує стандартизований, швидкий і надійний доступ до актуальних даних, 
дозволяє автоматизувати їх отримання та легко інтегрується з різними програмними системами. 
Ви запитуєте REST API подібно до стандартних HTTP-запитів, але майже завжди потрібно 
включати параметри. Параметри можуть передаватися як частина URL-адреси, через рядок 
запиту (Query Parameters), у заголовках HTTP або в тілі запиту залежно від 
використовуваного HTTP-методу.

Автентифікація API — це процес підтвердження особи користувача або програми перед 
наданням доступу до ресурсів сервера. Найпоширенішими способами є API Key, Bearer 
Token (JWT), Basic Authentication та OAuth 2.0.

Отримайте власний токен доступу на https://github.com/settings/tokens/new GitHub API 
використовує GET/PUT/DELETE, щоб ви могли автоматично запитувати або оновлювати 
елементи у вашому обліковому записі GitHub. 

\ding{43} Сервер не запам'ятовує 
ваші останні запити, ви завжди повинні включати свій токен доступу, 
якщо використовуєте його таким чином.

Приклад REST запиту:

\begin{lstlisting}
import os
from dotenv import load_dotenv

load_dotenv()
token = os.getenv("GITHUB_TOKEN")

headers = {'Authorization': 'Bearer '+token}
response = requests.get("https://api.github.com/user", headers=headers)
print(response.content)
\end{lstlisting}

Результат:

\begin{console}
b'{"login":"maxvonlancaster","id":38358502,"node_id":"MDQ6VXNlcjM4MzU4NTAy","avatar_url": ...
\end{console}

REST-архітектура базується на взаємодії клієнта і сервера, використанні стандартних HTTP-методів, 
відсутності збереження стану між запитами (stateless), єдиному інтерфейсі доступу до ресурсів та 
представленні ресурсів за допомогою URI. Дані найчастіше передаються у форматах JSON або XML.

HTTP-заголовки використовуються для передачі службової інформації між клієнтом і сервером. 
Вони можуть містити тип передаваних даних, параметри автентифікації, інформацію про кешування, 
мову, кодування, тип клієнта та інші параметри, необхідні для коректної обробки запиту або відповіді.



\section{CSV-файли.}

CSV (Comma-Separated Values): Відноситься до будь-якого текстового файлу з роздільниками 
(не завжди розділеного комами).

Якщо самі значення містять коми, ви можете взяти їх у лапки (наш реєстратор, очевидно, завжди 
так робить, щоб бути впевненим).

Імпортуємо дані із csv-файлу, що містить дані про історичну ціну біткоїна, та зобразимо ці 
дані, за допомогою наступного коду.

\begin{lstlisting}
import pandas as pd

dataframe = pd.read_csv("../resources/bitcoin.csv", delimiter=',', quotechar='"')
dataframe.head(10)
\end{lstlisting}

Результат:

\begin{console}
Date	Open	High	Low	Close	Volume	Currency
0	2010-07-18	0.0	0.1	0.1	0.1	75	    USD
1	2010-07-19	0.1	0.1	0.1	0.1	574	    USD
2	2010-07-20	0.1	0.1	0.1	0.1	262	    USD
3	2010-07-21	0.1	0.1	0.1	0.1	575	    USD
4	2010-07-22	0.1	0.1	0.1	0.1	2160	USD
5	2010-07-23	0.1	0.1	0.1	0.1	2403	USD
6	2010-07-24	0.1	0.1	0.1	0.1	496	    USD
7	2010-07-25	0.1	0.1	0.1	0.1	1551	USD
8	2010-07-26	0.1	0.1	0.1	0.1	877	    USD
9	2010-07-27	0.1	0.1	0.1	0.1	3374	USD
\end{console}

Формат CSV є простим, компактним і зрозумілим, підтримується більшістю програм для роботи з даними 
та легко обробляється різними мовами програмування. Він забезпечує швидке зчитування й запис 
табличних даних, що робить його одним із найпоширеніших форматів для обміну інформацією.

\ding{43} Формат CSV не підтримує вкладені структури даних, типи даних, метадані та форматування. 
Крім того, можуть виникати проблеми через використання різних роздільників, кодувань 
символів або символів нового рядка, що іноді ускладнює обмін даними між різними програмами.





\section{Парсинг JSON в Python.}

JSON (JavaScript Object Notation) — це легкий текстовий формат обміну даними, який використовується 
для передачі структурованої інформації між програмами. Незважаючи на назву, JSON сьогодні є 
незалежним від JavaScript і підтримується практично всіма сучасними мовами програмування.

JSON широко використовується у REST API, веб-сервісах, мобільних застосунках та хмарних системах. 
Основними перевагами JSON є компактність, простота читання людиною та легкість автоматичної обробки.

Вбудована бібліотека для читання/запису об'єктів Python з/у файли JSON, імпортується за допомогою 
команди `import json`. Нижче наведено приклад роботи із json-даними в python. REST-запит вертає 
json дані, які ми потім відображаємо.

\begin{lstlisting}
import json
# load json response from REST API

headers = {'Authorization': 'token ' + token}
response = requests.get("https://api.github.com/user", headers=headers)
data = json.loads(response.content)

#json.load(file) # load json from file
json.dumps(data) # return json string
#json.dump(obj, file) # write json to file
\end{lstlisting}

Результат:

\begin{console}
'{"login": "maxvonlancaster", "id": 38358502, "node_id": "MDQ6VXNlcjM4MzU4NTAy", "avatar_url": ...
\end{console}

JSON підтримує рядки, числа, логічні значення, масиви, об'єкти, а також спеціальне значення 
`null`. JSON-файл можна зчитати за допомогою стандартного модуля `json`, використовуючи 
функцію `json.load()`, або за допомогою бібліотеки `pandas`, якщо дані мають табличну 
структуру. 

Серіалізація JSON — це процес перетворення об'єкта Python у JSON-рядок або JSON-файл для 
його збереження чи передачі. Десеріалізація JSON — це процес перетворення JSON-рядка або 
JSON-файлу у відповідні об'єкти Python.

Після десеріалізації json в python значення вкладених полів отримують шляхом послідовного 
звернення до вкладених словників або списків за допомогою ключів та індексів.

Нижче наведено приклад роботи з вкладеними структурами в JSON:


\begin{lstlisting}
import json

json_string = '''
{
    "employee":
    {
        "name":"John",
        "department":
        {
            "name":"IT",
            "floor":4
        }
    }
}
'''

obj = json.loads(json_string)

print(obj["employee"]["department"]["name"])
\end{lstlisting}

Результат:

\begin{console}
IT
\end{console}




\section{Парсинг XML/HTML в Python.}

XML (eXtensible Markup Language) — це мова розмітки, призначена для зберігання, структурування та 
обміну даними між різними інформаційними системами. XML-документ складається з декларації, 
кореневого елемента, вкладених тегів, атрибутів та текстового вмісту. XML використовує теги 
для опису структури даних і є більш громіздким, тоді як JSON має компактніший синтаксис, 
легше читається людиною та швидше обробляється програмами. 

Існує ряд XML/HTML-парсерів для Python, але для науки про дані добре підходить бібліотека 
BeautifulSoup (спеціально призначена для вилучення даних з XML/HTML-файлів). Нижче наведено 
приклад завантаження html-розмітки із сайту, та код, що знаходить всі anchor-посилання 
всередені елементів `<div>`.

\begin{lstlisting}
from bs4 import BeautifulSoup
import requests

response = requests.get("https://mathmod.chnu.edu.ua/")

root = BeautifulSoup(response.content)
root.find("div").findAll("a")


item = root.find(class_="flex-item")
print(item.text)


titles = root.select("div.flex-item h2")

for t in titles:
    print(t.text)


link = root.find("a")

print(link["href"])
print(link.text)


table = root.find("table")

rows = table.find_all("tr")

for row in rows:
    cells = row.find_all("td")
    print([cell.text.strip() for cell in cells])
\end{lstlisting}

В цьому коді ми спочатку: отримуємо всі посилання на веб-сторінці, потім шукаємо 
елемент за класом, CSS селектором, отримуємо атрибути тега та, в кінці, отримуємо всі рядки таблиці.

Результат:

\begin{console}
[<a href="/cdn-cgi/l/email-protection#77141b105a1a16031f37141f190259121302590216"><span><span class="__cf_email__" data-cfemail="a0c3ccc78dcdc1d4c8e0c3c8ced58ec5c4d58ed5c1">[email protected]</span></span></a>,
 <a href="https://www.facebook.com/fmi.org.ua/" rel="noopener" target="_blank" title=".."><span class="iconify-icon footer-social-link-icon icon--fa-brands icon--fa-brands--facebook-f"></span></a>,
 <a href="https://www.instagram.com/m._i_.f/" rel="noopener" target="_blank" title=".."><span class="iconify-icon footer-social-link-icon icon--fa-brands icon--fa-brands--instagram"></span></a>,
 <a href="https://youtube.com/@MathTube_" rel="noopener" target="_blank" title="Youtube"><span class="iconify-icon footer-social-link-icon icon--fa-brands icon--fa-brands--youtube"></span></a>]
\end{console}

HTML-документ — це текстовий файл, який містить структуру веб-сторінки, описану за допомогою 
HTML-тегів. HTML є основним джерелом інформації під час веб-скрапінгу, оскільки саме в ньому 
містяться текст, таблиці, посилання, зображення та інші елементи веб-сторінки.

HTML-елемент за тегом можна знайти за допомогою методів 
\verb|find()| або \verb|find_all()|, вказавши 
назву потрібного тегу. Елементи знаходять, вказуючи значення атрибутів, наприклад 
"class", "id" або інших, використовуючи параметри методів \verb|find()| чи 
\verb|find_all()|.

DOM (Document Object Model) — це деревоподібне представлення HTML-документа, у якому 
кожен HTML-елемент є окремим вузлом, доступним для аналізу та модифікації.





\section{Регулярні вирази.}

Після завантаження даних (або якщо вам потрібно створити парсер для завантаження даних 
іншого формату) часто доводиться шукати певні елементи в даних.

Регулярні вирази (Regular Expressions, скорочено RegEx або Regex) — це спеціальна мова 
шаблонів, призначена для пошуку, перевірки, виділення та заміни текстових фрагментів. 
Регулярні вирази широко застосовуються під час збору, очищення та попередньої обробки 
даних, оскільки дозволяють швидко знаходити рядки, що відповідають заданому шаблону. 
У задачах інтелектуального аналізу даних регулярні вирази використовуються для пошуку 
даних (електронна пошта, номери телефонів, дати/час), перевірки правильності введених 
даних, очищення тексту від службових символів, та попередньої підготовки текстів перед 
їх аналізом.

Наприклад, знайти перше входження рядка «data science»

\begin{lstlisting}
import re
text = "This course will introduce the basics of data science"
match = re.search(r"data science", text)
print(match.start())
print(re.match(r"This", text))

match = re.match(r"data science", text) 
match = re.search(r"data science", text) 

all_matches = re.findall(r"a", text) 
print(all_matches)
\end{lstlisting}

Тут метод \verb|re.findall(r"a", text)| повертає всі збіги в тексті.

Результат:

\begin{console}
41
<re.Match object; span=(0, 4), match='This'>
['a', 'a', 'a']
\end{console}

Під час збору та підготовки даних регулярні вирази використовуються для автоматичного вилучення 
структурованої інформації з текстів. Типовими прикладами є пошук електронних адрес, телефонних 
номерів, URL-адрес, IP-адрес, поштових індексів, дат, часу, грошових сум та інших шаблонних даних.

Крім пошуку інформації, регулярні вирази широко застосовуються для очищення текстів перед аналізом: 
видалення HTML-тегів, спеціальних символів, зайвих пробілів, службових позначок та інших небажаних 
елементів. Вони також використовуються для перевірки правильності введених даних (валідації) та 
підготовки текстових корпусів до подальшої обробки методами обробки природної мови (Natural 
Language Processing, NLP).

Завдяки простоті використання та високій швидкодії регулярні вирази є одним із базових 
інструментів інженера даних і спеціаліста з інтелектуального аналізу даних.

\subsection{Відповідність декількох потенційних символів.}

Справжня сила регулярних виразів полягає в можливості збігу декількох можливих послідовностей 
символів. Спеціальні символи в регулярних виразах: \verb|.^$*+?{}\[]|()| (якщо 
ви хочете, щоб ці символи збігалися точно, вам потрібно їх екранувати: \verb|\$|)

Відповідність наборів символів:

\begin{itemize}
\item Відповідність символу «a»: \verb|a|
\item Відповідність символів «a», «b» або «c»: \verb|[abc]|
\item Будь-який символ, крім «a», «b» або «c»: \verb|[\^$abc]|
\item Відповідність будь-якої цифри: \verb|\d| (те саме, що `[0-9]`)
\item Відповідність будь-якого алфавітно-цифрового символу: \verb|\w| (те саме, що \verb|[a-zA-z0-9_]|)
\item Відповідність пробілу: \verb|\s| (те саме, що \verb|[ \t\n\r\f\v]|)
\item Відповідність будь-якого символу: `.` (включно з новою лінією з re.DOTALL)
\end{itemize}

Може збігатися з одним або декількома випадками символу (або набору символів)

Деякі поширені модифікатори:

\begin{itemize}
\item Відповідність символу «a» рівно один раз: `a`
\item Відповідність символу «a» нуль або один раз: `a?`
\item Відповідність символу «a» нуль або більше разів: `a*`
\item Відповідність символу «a» один або більше разів: `a+`
\item Відповідність символу «a» рівно n разів: `a{n}`
\end{itemize}

Можна поєднувати їх із збігом декількох символів:

- Відповідність усіх випадків «`<something>` science», де `<something>` — це алфавітно-цифровий рядок, що містить принаймні один символ:
\verb|`\w+\s+science`|

\subsection{Групування.}

Часто нам потрібно отримати більше інформації, ніж просто те, чи знайшли ми збіг чи ні (наприклад, ми можемо хотіти знати, який текст збігся).

Групування: оберніть частини регулярного виразу в дужки, щоб «запам'ятати» ці частини збігу 

\verb|`(\w+)\s([Aa]nalysis)`|.

\begin{lstlisting}
text = "This course will introduce the basics of data analysis"
match = re.search(r"(\w+)\s([Aa]nalysis)", text)
print(match.start(), match.groups())
\end{lstlisting}

Чому перед рядком стоїть «r»? Це дозволяє уникнути необхідності подвійного екранування рядків

Результат:

\begin{console}
41 ('data', 'analysis')
\end{console}


\subsection{Заміни.}

Регулярні вирази надають механізм для заміни одного тексту іншим текстом.

\begin{lstlisting}
better_text = re.sub(r"data analysis", r"schmada analysis", text)

better_text = re.sub(r"(\w+)\s([Aa])nalysis", r"\1 \2shnalysis", text)
print(better_text)
\end{lstlisting}

Щоб включити текст, який було збережено під час пошуку за допомогою груп, використовуйте екрановані послідовності 
\verb|\1|, \verb|\2|, … у тексті заміни

Результат:

\begin{console}
This course will introduce the basics of data ashnalysis
\end{console}


\subsection{Порядок і жадібне зіставлення.}

У регулярних виразах існує порядок операцій. `abc|def` відповідає рядкам «abc» або «def», а не 
«ab(c або d)ef». Ви можете обійти це, використовуючи дужки, наприклад `a(bc|de)f`. Це 
також створює групу, використовуйте `a(?:bc|de)f`, якщо ви не хочете її захоплювати.

За замовчуванням регулярні вирази намагаються захопити якомога більше тексту 
(жадібне співставлення, greedy matching). `<(.*)>`, застосоване до `<a>text</a>`, буде відповідати всьому виразу. 
Якщо ви хочете захопити якомога менше тексту, використовуйте `<(.*?)>`, це буде відповідати тільки 
терміну `<a>`.





\medskip
\ding{46} \textbf{Питання для самоконтролю.}

\begin{enumerate}
\item Що таке збір даних (Data Collection) та яке його значення в процесі інтелектуального аналізу даних?
\item Які основні джерела даних використовуються в Data Mining?
\item Які переваги має автоматизований збір даних порівняно з ручним?
\item Які труднощі можуть виникати під час збору даних з різних джерел?
\item Що таке веб-джерело даних?
\item Які типи веб-джерел можуть використовуватися для збору даних?
\item Що таке веб-скрапінг (Web Scraping)?
\item Які етапи процесу веб-скрапінгу?
\item Які бібліотеки Python найчастіше використовуються для отримання веб-сторінок?
\item Для чого використовується бібліотека `requests`?
\item Що таке HTTP-запит?
\item Які основні HTTP-методи використовуються при роботі з веб-ресурсами?
\item Що таке HTTP-відповідь?
\item Яке призначення кодів стану HTTP?
\item Що означають коди стану 200, 404 та 500?
\item Що таке RESTful API?
\item Які принципи лежать в основі REST-архітектури?
\item Чим REST API відрізняється від звичайної веб-сторінки?
\item Які переваги використання REST API для збору даних?
\item Як передаються параметри в REST API-запитах?
\item Що таке JSON-відповідь REST API?
\item Для чого використовуються HTTP-заголовки (headers)?
\item Що таке автентифікація API?
\item Які способи автентифікації API ви знаєте?
\item Що таке CSV-файл?
\item Які особливості структури CSV-файлів?
\item Які переваги має формат CSV?
\item Які недоліки має формат CSV?
\item Як зчитати CSV-файл у Python?
\item Для чого використовується бібліотека `pandas` при роботі з CSV?
\item Що таке JSON?
\item Які типи даних підтримує формат JSON?
\item Як зчитати JSON-файл у Python?
\item Що таке серіалізація JSON?
\item Що таке десеріалізація JSON?
\item Як отримати значення вкладених полів у JSON-структурі?
\item Що таке XML?
\item Які основні елементи XML-документа?
\item Чим XML відрізняється від JSON?
\item Як виконати парсинг XML у Python?
\item Що таке HTML-документ?
\item Яку роль HTML відіграє у веб-скрапінгу?
\item Для чого використовується бібліотека BeautifulSoup?
\item Як знайти HTML-елемент за тегом?
\item Як знайти HTML-елемент за атрибутом або класом?
\item Що таке DOM-структура веб-сторінки?
\item Що таке регулярний вираз (Regular Expression)?
\item Для яких задач використовуються регулярні вирази під час збору даних?
\item Які основні метасимволи регулярних виразів ви знаєте?
\item Як за допомогою регулярних виразів знаходити та витягувати потрібні фрагменти тексту?
\end{enumerate}




